\chapter{习题与推导练习}
本附录的练习对应正文中的关键推导、近似条件和测量判据。建议先独立完成，再核对解题要点；用于研究计算时，可将解析极限编写为数值单元测试。题目按基本推导和模型检验组织，题目中的可复核结论均可在正文找到。

\section{基本推导与概念检验}
\subsection*{A1：从轨迹积分推出相位匹配带宽}
令 $E_z(z)=E_0\ee^{\ii k_zz}$ 只在 $-L/2<z<L/2$ 非零。直接计算 $\beta$，证明 $|\beta|^2$ 的主瓣宽度按 $1/L$ 缩小，并给出第一零点的 $\Delta k_z$。然后解释为什么“更长作用长度”同时意味着更强耦合和更窄的速度接受带宽。

\textit{解题要点。} 积分给出 $\beta\propto L\,\mathrm{sinc}[(k_z-\omega/v)L/2]$。第一零点满足 $|(k_z-\omega/v)L/2|=\pi$，故单侧零点宽度 $2\pi/L$。实验上若电子能散使 $v$ 分布跨过该宽度，继续增加 $L$ 不会按 $L^2$ 提高总边带概率。

\subsection*{A2：证明贝塞尔边带归一化}
从 $U_e(\beta)=\exp(-\beta b^\dagger+\beta^*b)$ 出发，证明 $P_n=J_n^2(2|\beta|)$ 并验证 $\sum_nP_n=1$。再给出一个数值截断准则，使遗漏概率低于 $10^{-8}$。

\textit{解题要点。} 在相位表象中取 $b^\dagger\to \ee^{-\ii\theta}$、$b\to \ee^{\ii\theta}$，并令 $\beta=|\beta|\ee^{\ii\phi}$，则 $U_e$ 变成周期因子 $\exp[-2\ii|\beta|\sin(\phi-\theta)]$，Fourier 系数即 Jacobi--Anger 展开。归一化来自幺正性或恒等式 $\sum_nJ_n^2(x)=1$。数值上应自适应增加 $n_{\max}$ 直到 $1-\sum_{|n|\le n_{\max}}J_n^2(2|\beta|)<10^{-8}$。

\subsection*{A3：反冲误差为什么随边带阶数放大}
从相对论能量 $E(p)=\sqrt{m^2c^4+p^2c^2}$ 在 $p_0$ 附近展开到二阶，重新推出前文的 $1/\gamma^3m$ 曲率，并估计给定 $\hbar\omega$ 时 $n$ 增大如何破坏等间隔能量晶格。

\textit{解题要点。} $\dd E/\dd p=v$，$\dd ^2E/\dd p^2=1/(\gamma^3m)$。线性项给 $\Delta p_n\simeq n\hbar\omega/v$，二阶项相对线性交换按 $|n|\hbar\omega/(2\gamma^3mv^2)$ 增长。真正的截断标准应在实验实际占据的 $n_{\max}$ 处检查，而不是只看单光子反冲。

\subsection*{A4：经典场极限不是简单令 $N\to\infty$}
对相干态 $|\alpha\rangle$，在 $U(g)=\exp(g a^\dagger b-g^*b^\dagger a)$ 的约定下保持 $\beta=g^*\alpha$ 固定，让 $|\alpha|\to\infty$。说明为什么必须同时令 $g\to0$，并估计相对光子数涨落如何消失。

\textit{解题要点。} 若 $g$ 不缩放，$\beta$ 会发散，比较的不是同一个经典问题。固定 $\beta$ 时 $g\sim1/|\alpha|$；相干态 $\Delta N/\bar N=1/\sqrt{\bar N}$，因此量子涨落修正随占据增加而下降。

\medskip
\subsection*{B1：什么时候二维能量--OAM 晶格真的二维}
给定两条跃迁向量 $\mathbf a_1=(1,\ell_1)$、$\mathbf a_2=(r,\ell_2)$，求它们线性无关的条件。若行列式为零，说明为什么再复杂的二维绘图也只是一个一维子格的重标记。

\textit{解题要点。} 条件为 $\det(\mathbf a_1,\mathbf a_2)=\ell_2-r\ell_1\neq0$。若为零，所有可达格点都落在同一直线生成的子群上，二维坐标只是冗余标签。

\subsection*{B2：Lifshitz 转变与 Chern 转变为什么不是一回事}
对单带色散 $\varepsilon(\vb{k})$ 求鞍点并讨论零能等高线重连；再解释为什么单个标量色散没有足够的本征矢结构定义非平凡 Berry 曲率。构造一个最小两带 $\mathbf d(\vb{k})\cdot\boldsymbol\sigma$ 模型作为对照。

\textit{解题要点。} Lifshitz 转变只需要费米面/等能线拓扑改变；Chern 数依赖本征矢在整个 BZ 的映射。标量单带模型的周期本征矢可以取常数，Berry 曲率为零。只有矩阵哈密顿量的 $\hat{\mathbf d}$ 发生非平凡缠绕时才有非零 Chern 数。

\subsection*{B3：推导 Jackiw--Rebbi 零模局域长度}
从\cref{eq:ch03-jr-zero} 出发，取 $\kappa(z)=\kappa_0\tanh(z/w)$，求零模包络，并讨论 $w\ll\ell_{\rm JR}$ 与 $w\gg\ell_{\rm JR}$ 两个极限。

\textit{解题要点。} 积分 $\int\tanh(z/w)\dd z=w\ln\cosh(z/w)$，故 $|\psi_0|\propto[\cosh(z/w)]^{-\kappa_0w/(\hbar v_D)}$。远离畴壁恢复指数衰减长度 $\hbar v_D/\kappa_0$；畴壁很宽时中心区域更接近高斯/缓变局域。

\subsection*{B4：二能级 Bragg 截断的 Schrieffer--Wolff 修正}
从三能级 $|{-1}\rangle,|0\rangle,|1\rangle$ 模型出发，令 $|-1\rangle$ 相对目标二能级失谐 $\Delta\gg g$。消去 $|-1\rangle$，求 $O(g^2/\Delta)$ 的 AC Stark 位移，并说明它如何造成第八章所说的系统相位偏差。

\textit{解题要点。} 二阶虚跃迁给目标态能量修正约 $\pm |g|^2/\Delta$（符号取决于失谐约定）。即使泄漏概率 $O[(g/\Delta)^2]$ 很小，相位在长作用时间中仍可积累到可测量尺度，所以“看不到第三边带”不等于二能级模型无偏。

\section{测量与模型检验}
\subsection*{C1：comb 态为什么是能量基的 Fourier 共轭}
从\cref{eq:ch04-comb-state} 证明 $b_\Omega$ 的本征关系，并对有限 $d$ 维循环能量格点写出归一化离散 comb 基。证明能量基和 comb 基满足 $|\langle n|\phi_j\rangle|^2=1/d$。

\textit{解题要点。} 有限维取 $|\phi_j\rangle=d^{-1/2}\sum_{n=0}^{d-1}e^{-2\pi\ii jn/d}|n\rangle$。任意能量态对任意 comb 态的重叠模方均为 $1/d$，所以两基互为 mutually unbiased bases。

\subsection*{C2：为什么只测电子会丢掉参数信息}
取一个参数化纯联合态 $|\Psi(\theta)\rangle=\cos\theta|00\rangle+\sin\theta|11\rangle$。比较联合计算基测量、只测第一个子系统、以及把第二个子系统完全迹掉后的 Fisher 信息。再构造一个相位参数 $|00\rangle+\ee^{\ii\theta}|11\rangle$，说明计算基人口为什么完全看不到相位。

\textit{解题要点。} 对振幅参数，联合人口和单边人口都能保留部分/全部经典信息；对纯相位参数，计算基人口与 $\theta$ 无关，Fisher 信息为零，但换到互补基或做联合干涉即可恢复。这正是测量基选择的核心。

\subsection*{C3：非破坏测量也会有回作用}
对 comb 态与辅助光的联合演化\cref{eq:ch04-comb-displacement}，考虑两个不同相位 $\phi_1,\phi_2$ 的电子叠加。把光场迹掉，证明电子非对角元被两个位移光态的重叠 $\langle\alpha_2|\alpha_1\rangle$ 乘上。解释“读出越可区分，电子相位越被投影”这一测量--退相干互补性。

\textit{解题要点。} 若辅助光初态为真空，输出为 $|\alpha_j\rangle$，偏迹后 $\rho_{12}\to\rho_{12}\exp[-|\alpha_1-\alpha_2|^2/2]\ee^{\ii\Phi}$。所谓非破坏指不必吸收/销毁电子，不代表没有量子回作用。

\subsection*{C4：用 Fisher 矩阵设计两种互补测量}
设未知参数为 $(|\beta|,\phi)$。构造只测边带人口的前向模型，再加入一个已知参考耦合。证明前者在弱耦合下对绝对相位的 Fisher 信息为零，而参考相位扫描可以使 Fisher 矩阵满秩。

\textit{解题要点。} 单次 $J_n^2(2|\beta|)$ 不含 $\arg\beta$；加入 $\beta_r \ee^{\ii\phi_r}$ 后概率依赖 $|\beta+\beta_r \ee^{\ii\phi_r}|$。至少两个非共线参考相位才能同时约束复平面的两个方向。

\medskip
\subsection*{D1：Bragg--PINEM 交叉区}
建立至少五个能量格点的含反冲模型，扫描 $\mathcal R=|g|/\Delta_{\rm rec}$。分别用完整模型和二能级模型拟合 Rabi 频率与相位。找出“第三边带仍低于 1\%，但二能级相位偏差已经超过 0.1 rad”的参数区。如果不存在，就否决“隐形多通道偏差”作为近期实验课题。

\subsection*{D2：一阶--二阶路径干涉}
对含有 $g_1|n+1\rangle\langle n|+g_2|n+2\rangle\langle n|+\mathrm{h.c.}$ 的能量阶梯做时间演化，扫描 $\Phi=\phi_2-2\phi_1$。加入能量分辨率卷积和相位噪声，计算区分 $g_2=0$ 与 $g_2\neq0$ 所需电子数。若所需剂量超过样品可承受剂量两个数量级，则不要把理想相位条纹写成“可实验验证”。

\subsection*{D3：comb-basis 闭环纠错}
模拟宽度 $\sigma_k$ 的高斯能量梳经过随机二次相位误差 $\delta$，用有限压缩辅助光产生 noisy comb readout，再根据后验均值反馈。比较无反馈、一次反馈和连续弱反馈的逻辑保真度。必须同时计入辅助测量造成的退相干；若反馈提升来自不公平地忽略测量回作用，结果无效。

\subsection*{D4：低能拓扑零模的可观测性}
用 $\kappa(z)=\kappa_0\tanh(z/w)$ 的 Dirac 模型生成零模，再恢复共同的二次色散项和有限能量宽度。设计一个时间飞行或能谱观测量，比较拓扑畴壁与一个经过参数拟合的普通势阱。只有找到普通势阱无法同时模仿的特征，才能把观测称为拓扑证据。

\subsection*{D5：自由电子读取发射体相干性的最优探针}
以\cref{eq:ch07-coherence-interference} 为最小前向模型，把电子入射相位、延迟和能量窗口作为设计变量，目标是最大化对 $\operatorname{Re}\rho_{ge}$ 与 $\operatorname{Im}\rho_{ge}$ 的最小 Fisher 本征值。加入未知人口和背景作为 nuisance parameters。若最优设计仍出现近奇异 Fisher 矩阵，就说明仅靠该电子谱不能完成完整相干层析，需要第二种测量基。

\subsection*{D6：能量--OAM 后选择的多模量子光 Kraus 工程}
构造两个具有不同 OAM 的量子光模式，与电子产生两条在能量上简并、在 OAM 上可区分的交换路径。推导联合后选择 $M_{n\ell}$，比较只测能量的 $M_n$。目标不是让条件 Wigner 函数“更漂亮”，而是证明 OAM 分辨使某个原本秩亏的光子态参数反演变为满秩，并计算成功概率损失。

\medskip
每道研究题建议固定写四栏：模型与近似、无量纲控制参数、旧模型严格禁止的观测量、否决条件。只有前三栏没有第四栏，很容易把任何复杂数值结果解释成“新效应”。真正能长成论文的题目，通常允许你在计算早期就说清楚什么结果会让自己放弃这个想法。

把解析推导变成代码时，还应至少保留三个自动测试：零耦合回到自由传播；关闭新增项回到上一层模型；扩大数值截断后结果收敛。若一个“新现象”随着边带截断、网格或时间步改变而漂移，它首先是数值问题，不是物理发现。
